% TOP_PROBLEM: {"problemId": "problem.vassiliev-finite-type-invariant-separation-conjecture", "canonicalTitle": "Vassiliev Separation by Abelian-Valued Finite-Type Invariants", "releaseRank": 288, "primaryDomain": "geometry_topology", "primaryDomainLabel": "Geometry & topology", "displayStatus": "open", "releaseStatus": "open"}
% !TeX program = lualatex
% Definition and sources only; this file is not an LLM solution attempt.
\documentclass[11pt]{article}
\usepackage[margin=25mm]{geometry}
\usepackage{fontspec}
\setmainfont{Latin Modern Roman}
\usepackage{amsmath,amssymb,mathtools}
\usepackage{unicode-math}
\setmathfont{Latin Modern Math}
\usepackage{xurl,hyperref}
\hypersetup{colorlinks=true,urlcolor=blue}
\setcounter{secnumdepth}{0}
\setlength{\parindent}{0pt}
\setlength{\parskip}{5pt}
\setlength{\emergencystretch}{3em}
\begin{document}
\section*{\texorpdfstring{Vassiliev Separation by Abelian-Valued Finite-Type Invariants}{Vassiliev Separation by Abelian-Valued Finite-Type Invariants}}\label{288}
\subsection{Definitions and mathematical statement}
Let $\mathcal K$ be the set of smooth closed oriented unframed knots in $S^3$ modulo orientation-preserving ambient isotopy. An invariant $v:\mathcal K\to A$, with $A$ abelian, extends linearly to ${\mathbb{Z}}[\mathcal K]$ and to singular knots using
$v(K_\times)=v(K_+)-v(K_-)$. It has finite type at most $d$ if it vanishes on every singular knot with $d+1$ transverse double points. The target is
\[
 K_1\not\cong K_2\quad\Longrightarrow\quad
 \exists A,d,v\text{ of this type with }v(K_1)\ne v(K_2).
\]
The coefficient group and order may depend on the pair. A knot and its reverse are included if they are distinct oriented types. Rational-valued separation is not silently substituted for the arbitrary-abelian-coefficient version.
\par\smallskip\textit{Formulation and concept references:} \href{https://doi.org/10.2140/agt.2026.26.1037}{[S1]}.

\subsection{Short English statement}
Can any two distinct oriented knot types be distinguished by some finite-type invariant, allowing its finite order and abelian coefficient group to depend on the pair?
\subsection{Sources}
\begin{itemize}
\item[S1]Bruno Dular, Primitive Feynman Diagrams and the Rational Goussarov–Habiro Lie Algebra of String Links, AGT 26:3 (2026), pp. 1037–1038, definitions and m=1 separation question; Ohtsuki problem list for oriented-knot formulation.
\url{https://doi.org/10.2140/agt.2026.26.1037}.
\textit{Formulation source.}
\end{itemize}
\end{document}
